% cg-layout: comparison
% Where a citekey is checked: a typical retrieval-augmented writer beside
% this pipeline, part by part. The rows carry the comparison: three of
% the four parts are the same kind of thing in both panels, so the eye
% runs across each row and stops at the one where (a) has nothing.
%
% WIDTH. Laid out to the render's \textwidth (453pt, A4 at 12pt with 1in
% margins) and \input bare. Two panels side by side take about three
% fifths of it; stacked, the same panels would take twice the height for
% no gain, which is the sprawl this layout exists to prevent.
% >>> chitragupta figure style v5 -- copy of assets/tikz/cg-figstyle.tex, do not edit <<<
% ===========================================================================
% chitragupta house figure style -- the "presentation tier" prelude block.
%
% WHAT THIS IS. The five \definecolor lines of docs/TIKZ-STYLE.md, grown
% into the rest of the vocabulary a figure needs to look drawn rather
% than assembled: zone cards, two-tier node labels, step badges, edge
% tags, and five domain `pic`s. Same travel rule as the palette -- this
% block lives *inside* the figure file, not in a shared .sty and not in a
% preamble, because the renderer injects only \usepackage{tikz} and a
% thesis-chapter fragment is \input into a document that has never heard
% of this project.
%
% WHY IT IS SAFE TO REPEAT. Every definition below is \providecommand or
% a \tikzset key. \tikzset is idempotent -- re-running it re-defines the
% same keys to the same values -- so N figures in one document carrying N
% copies of this block cost N key definitions and change nothing. That is
% the same property \usetikzlibrary has under #781, for the same reason.
% It is NOT the property \tikzset{every node/.style=...} has, which is
% why nothing here touches `every node`.
%
% LIBRARIES. Load these at the top of the figure, above \begin{tikzpicture}:
%   \usetikzlibrary{arrows.meta,positioning,fit,backgrounds,calc,shadows.blur}
% All six are TikZ/PGF libraries shipped by texlive-pictures. `shadows.blur`
% comes from pgf-blur and is reached by \usetikzlibrary, not \usepackage --
% verified, which is what makes the soft shadow portable into a foreign
% thesis.
%
% ONE TYPE SIZE: THE DOCUMENT'S OWN. Every word in a figure is set at
% \normalsize, the body size of the paragraph beside it -- glosses, edge
% tags, notes, legend text, badge digits, node names and zone titles
% alike. Roles are told apart by weight and colour, not size: a name or a
% zone title is bold, a gloss is muted.
%
% A figure whose labels run smaller than the paragraph beside them reads
% as an afterthought, and shrinking type is never the fix for a box with
% too many words in it -- that is the conciseness defect, and the fix is
% fewer words. Larger is not free either (v4): v3 set names and titles
% at \large, and a figure then read louder than its text and a title
% grew past its zone into the next one.
%
% DRAW AT FINAL WIDTH. \resizebox NEVER. This is the same rule from the
% other end, and it is the one that actually gets broken, because a
% \resizebox is invisible in the figure file -- it lives in the draft.
% A picture set at its natural size and scaled to fit prints its type at
% whatever the scale factor happens to be, so the floor above becomes a
% number nobody chose. Measured on one acmart paper before the fix: six
% figures scaled between 0.92 and 1.18, with the smallest type landing
% between 5.5pt and 7.1pt against a 9pt body (#1012).
%
% So: settle the target width first (docs/TIKZ-STYLE.md, "Draw for the
% width it will be printed at"), lay the figure out to that width, and
% \input it bare. If it does not fit, change the layout -- wrap a row,
% cut a gloss, drop a stage -- never the scale.
% ===========================================================================

% --- Palette --------------------------------------------------------------
% The five roles of docs/TIKZ-STYLE.md, unchanged, plus three more members
% of the same Okabe-Ito set. The extension is not a new palette: a zoned
% figure needs one hue per zone and five roles do not stretch to six zones
% without reusing one, which reads as a claim the figure is not making.
\definecolor{cgInk}{HTML}{1A1A1A}%      borders, labels, default strokes
\definecolor{cgFlow}{HTML}{0072B2}%     the primary path through the figure
\definecolor{cgAccent}{HTML}{D55E00}%   the one thing the figure is about
\definecolor{cgAlt}{HTML}{009E73}%      a second class of node or edge
\definecolor{cgAux}{HTML}{56B4E9}%      auxiliary, "also happens"
\definecolor{cgWarm}{HTML}{E69F00}%     a third zone   (Okabe-Ito orange)
\definecolor{cgRose}{HTML}{CC79A7}%     a fourth zone  (Okabe-Ito purple)
\definecolor{cgPaper}{HTML}{FAFAF8}%    card ground, one step off white
% A named grey rather than `cgInk!55`. Styles here take a colour as an
% argument and paste it into a tint expression (`fill=#1!11`), and
% `cgInk!55!11` is not a colour -- xcolor stops with "Undefined color".
% Every role that needs muted ink therefore needs a *name*.
\definecolor{cgMute}{HTML}{8A8A8A}%     "absent", and named as such
%
% --- Derived shades: a palette name fills; a Line or Text name draws -------
% The hues above are fills. As strokes several are too light to see on
% their own tint (cgWarm's border on cgWarm!12 is 1.9:1, WCAG asks 3:1
% of a line), and as text none but cgFlow reaches the 4.5:1 a label
% needs. So each hue has two darker shades, mixed towards cgInk:
% `cg<Hue>Line` for any border, edge or rule, `cg<Hue>Text` for any
% coloured text. Each share is the largest multiple of 5 that clears the
% threshold over white, cgPaper and every tint this block fills with;
% tests/test_figure_contrast.py recomputes it from the hex values above
% rather than trusting these numbers. A Text shade clears 4.5:1 against
% white, so white type on a Text fill (the step badge) clears it too.
\colorlet{cgFlowLine}{cgFlow!100!cgInk}\colorlet{cgFlowText}{cgFlow!95!cgInk}%
\colorlet{cgAccentLine}{cgAccent!100!cgInk}\colorlet{cgAccentText}{cgAccent!75!cgInk}%
\colorlet{cgAltLine}{cgAlt!95!cgInk}\colorlet{cgAltText}{cgAlt!70!cgInk}%
\colorlet{cgAuxLine}{cgAux!75!cgInk}\colorlet{cgAuxText}{cgAux!55!cgInk}%
\colorlet{cgWarmLine}{cgWarm!75!cgInk}\colorlet{cgWarmText}{cgWarm!55!cgInk}%
\colorlet{cgRoseLine}{cgRose!90!cgInk}\colorlet{cgRoseText}{cgRose!65!cgInk}%
\colorlet{cgMuteLine}{cgMute!95!cgInk}\colorlet{cgMuteText}{cgMute!70!cgInk}%

\tikzset{
  % --- Picture-wide defaults ----------------------------------------------
  % Set once as a picture option: \begin{tikzpicture}[cg]
  % `font=\normalsize` is load-bearing rather than a default: every word
  % in a figure is set at the document's body size. Set it explicitly
  % rather than inheriting it, so a picture that lands inside a float,
  % a table cell or a \small group does not quietly drop below it.
  %
  % This only holds if the figure is drawn at its final width. A picture
  % set at natural size and then wrapped in \resizebox prints its type at
  % whatever the scale factor happens to be -- see the FINAL WIDTH note
  % at the head of this file.
  %
  % `node distance=6mm` is the gap between two members of one zone. Two
  % zones sit 8mm apart card to card, which with `cgzone`'s 3mm inner sep
  % is 14mm member to member: say that distance where you place across.
  cg/.style={
    thick, line cap=round, line join=round,
    font=\normalsize, node distance=6mm,
    every node/.append style={text=cgInk},
  },
  % --- Soft shadow --------------------------------------------------------
  % The single largest contributor to "drawn, not assembled". Kept very
  % light: a shadow that reads as a shadow is a 2008 slide deck.
  cgsoft/.style={
    blur shadow={shadow blur steps=10, shadow xshift=0.35pt,
                 shadow yshift=-0.7pt, shadow opacity=22,
                 shadow blur radius=1.6pt}},
  % --- Nodes --------------------------------------------------------------
  % One base box, four role tints. Rounded corners at 3pt: enough to read
  % as a card, not so much it reads as a pill.
  cgbox/.style={draw=cgFlowLine, fill=cgFlow!7, rounded corners=3pt,
                align=center, inner sep=4.5pt, minimum height=11mm,
                text width=24mm, cgsoft},
  cgalt/.style={cgbox, draw=cgAltLine,    fill=cgAlt!9},
  cgaux/.style={cgbox, draw=cgAuxLine,    fill=cgAux!10},
  cgwarm/.style={cgbox, draw=cgWarmLine,  fill=cgWarm!12},
  cgrose/.style={cgbox, draw=cgRoseLine,  fill=cgRose!10},
  cgplain/.style={cgbox, draw=cgMuteLine, fill=cgPaper},
  % The accent node. Heavier rule, because one thing per figure is what
  % the figure is about and the reader should find it without reading.
  cgkey/.style={cgbox, draw=cgAccentLine, fill=cgAccent!11, line width=1.15pt},
  % "Absent / not done here" -- dashed, unfilled, same footprint. A gap in
  % a pipeline is a claim and should occupy space like one.
  cgnone/.style={cgbox, draw=cgMuteLine, fill=white, cgflat,
                 dash pattern=on 2.6pt off 2.2pt},
  % Optional / off to the side: dashed, tinted, and deliberately *not*
  % shadowed. A dashed border plus a soft shadow renders as two parallel
  % broken lines, which reads as a drawing mistake rather than a dash.
  cgopt/.style={cgbox, draw=cgAuxLine, fill=cgAux!8, cgflat,
                dash pattern=on 2.6pt off 2.2pt},
  % Cancels the shadow a role style already set.
  cgflat/.style={blur shadow={shadow opacity=0, shadow blur steps=1}},
  % An icon well: an empty, invisible node of fixed size that a `pic`
  % is drawn into. It exists so the icon has a name -- which is what
  % lets a zone card `fit` around it and what lets `review figure`
  % measure it, since the aid sees only nodes the source names.
  cgwell/.style={inner sep=0pt, minimum width=10mm, minimum height=10mm},
  % --- Zone card ----------------------------------------------------------
  % The bRAG idiom: a tinted rounded region drawn *behind* its members,
  % with its title outside the top-left corner in the zone's own hue.
  % Draw it on the `background` layer so it never paints over a member.
  % The 3mm inner sep and the title's 1.5mm lift are as tight as a card
  % goes and still reads as a region rather than a border round a box.
  % Takes a palette name (`cgzone=cgAlt`), and draws in its Line shade.
  cgzone/.style={rounded corners=5pt, inner sep=3mm, draw=#1Line, fill=#1!6,
                 dash pattern=on 3.4pt off 2.6pt, line width=0.9pt},
  cgzone/.default=cgFlow,
  cgzonetitle/.style={anchor=south west, font=\normalsize\bfseries, text=#1Text,
                      inner sep=0pt, yshift=1.5mm},
  cgzonetitle/.default=cgFlow,
  % --- Edges --------------------------------------------------------------
  cgedge/.style={draw=cgFlowLine, line width=0.85pt,
                 -{Latex[length=2.6mm, width=2.1mm]}, shorten >=1.5pt,
                 shorten <=1pt},
  cgweak/.style={cgedge, draw=cgAuxLine, dash pattern=on 2.8pt off 2.2pt},
  cgback/.style={cgedge, draw=cgAccentLine, dash pattern=on 2.8pt off 2.2pt},
  % A leader line from a zone card to the point on the spine it explains.
  % No arrowhead: it points at a place, it does not carry anything there.
  % No opacity: a translucent Line shade is a lighter one, and the Line
  % shade is already the lightest that clears 3:1.
  cglead/.style={draw=#1Line, line width=0.9pt,
                 dash pattern=on 2.2pt off 2.4pt},
  cglead/.default=cgFlow,
  % --- Edge tag -----------------------------------------------------------
  % A word beside an arrow, never on it. A tag filled over its line
  % breaks the line where the eye is following it, so `auto` sets it on
  % the left of the direction of travel (`swap` for the right) with about
  % 1mm between the line and the type: near enough to belong to the
  % segment, far enough not to touch it.
  cgtag/.style={font=\normalsize, auto, inner sep=0.5mm, outer sep=0.5mm,
                text=cgMuteText},
  cgnote/.style={font=\normalsize\itshape, text=cgMuteText, inner sep=2pt},
  % --- Step badge ---------------------------------------------------------
  % A numbered disc at a node's top-left corner. Ordinal information
  % leaves the label, so the label is free to say what the step does.
  % At the type floor a two-digit badge needs 5.8mm, not 4.6mm. The disc
  % is sized from the type rather than the type shrunk to the disc. The
  % disc is the Text shade, which is what gives the white digit 4.5:1.
  cgbadge/.style={circle, draw=#1Text, fill=#1Text, text=white,
                  font=\normalsize\bfseries, inner sep=0pt,
                  minimum size=5.8mm},
  cgbadge/.default=cgFlow,
  % --- Legend -------------------------------------------------------------
  % A swatch that is the same object as the thing it stands for: the
  % card's rule weight and tint, its corner rounded, at 4mm square. A legend
  % drawn with a plain filled square claims a correspondence the figure
  % does not actually have.
  cgswatch/.style={draw=#1Line, fill=#1!11, rounded corners=1.2pt,
                   minimum size=4mm, inner sep=0pt, line width=0.8pt},
  cgswatch/.default=cgFlow,
  cglegendtext/.style={font=\normalsize, text=cgInk!85, anchor=west,
                       inner sep=0pt},
}%

% `\cglegenditem{<name>}{<anchor>}{<colour>}{<text>}` -- one legend entry,
% placed with its swatch at `<anchor>` and its text to the right.
\providecommand{\cglegenditem}[4]{%
  \node[cgswatch=#3] (#1) at (#2) {};
  \node[cglegendtext, right=1.6mm of #1] (#1-t) {#4};}%

% `\cgstep{<n>}{<node>}{<colour>}` -- the badge, placed. Written as a macro
% rather than a style because it needs the node's north-west corner and a
% style cannot reach for one.
\providecommand{\cgstep}[3]{%
  \node[cgbadge=#3] at ([xshift=1.2mm,yshift=-1.2mm]#2.north west) {#1};}%

% `\cgzonelabel{<node or fit box>}{<colour>}{<text>}` -- a zone's title.
\providecommand{\cgzonelabel}[3]{%
  \node[cgzonetitle=#2] at (#1.north west) {#3};}%

% --- The two-tier label ----------------------------------------------------
% `\cglab{<name>}{<gloss>}` -- a bold name over a muted gloss line.
%
% A name and a gloss are two roles, and setting them alike is what makes
% a box of text out of a label. They are told apart by weight and colour
% at one size, the body's: bold ink for the name, muted for the gloss.
% This is the only way a gloss gets set, which holds the roles apart
% across the figure. Keep the gloss to six words or fewer (v5; `figure
% check`'s w3-gloss); past that it is prose in a box and the conciseness
% defect is real again.
\providecommand{\cglab}[2]{%
  {\normalsize\bfseries #1}\\[1.5pt]{\normalsize\textcolor{cgMuteText}{#2}}}%

% ===========================================================================
% Domain `pic`s -- V.2's "a few house shapes, not an icon set".
%
% Each is drawn from explicit coordinates inside its own definition, which
% is the carve-out FIGURE-ROADMAP.md's "Rules that will have to bend"
% proposes: relative placement governs composition, a fixed object's
% internal geometry does not reflow and may be written down. Each is sized
% in millimetres against a body-text figure and each takes its colour from
% `cg pic colour`, so a pic inherits the zone it sits in: a palette name,
% whose Line shade draws the outline (`\cgPicColour Line`). The paler
% inner strokes are texture, not the shape, and stay tints. Each fits the
% 10mm `cgwell` it is drawn into; tests/test_tikz_scaffolds.py measures
% that.
%
% Each also passes the twin test: the shape is redundant with the label
% beside it. A drum labelled `ledger` and the twin's [ledger] carry the
% same proposition, so nothing is lost when the ASCII form drops the drum.
% ===========================================================================
\tikzset{
  cg pic colour/.store in=\cgPicColour, cg pic colour=cgFlow,
  % A drum: any persistent store. 8mm wide, 9mm tall. Twin: [ledger].
  pics/cgstore/.style={code={
    \path[fill=\cgPicColour!15, draw=\cgPicColour Line, line width=0.8pt]
      (-4mm,2.6mm) -- (-4mm,-3mm)
      .. controls (-4mm,-4.4mm) and (4mm,-4.4mm) .. (4mm,-3mm) -- (4mm,2.6mm)
      .. controls (4mm,4mm) and (-4mm,4mm) .. (-4mm,2.6mm) -- cycle;
    \path[draw=\cgPicColour Line, line width=0.8pt]
      (-4mm,2.6mm) .. controls (-4mm,1.2mm) and (4mm,1.2mm) .. (4mm,2.6mm);
  }},
  % A stack of pages: a document, or a set of them. Twin: [PDFs].
  pics/cgdocs/.style={code={
    \foreach \i in {2,1}{
      \path[fill=cgPaper, draw=\cgPicColour!55, line width=0.7pt,
            rounded corners=0.6pt]
        (-3.6mm+\i*0.8mm, -4.5mm+\i*0.8mm) rectangle (2mm+\i*0.8mm, 2.9mm+\i*0.8mm);}
    \path[fill=cgPaper, draw=\cgPicColour Line, line width=0.8pt, rounded corners=0.6pt]
      (-3.6mm,-4.5mm) rectangle (2mm,2.9mm);
    \foreach \y in {1.4,-0.1,-1.6,-3.1}{
      \path[draw=\cgPicColour!60, line width=0.55pt]
        (-2.6mm,\y mm) -- (1mm,\y mm);}
  }},
  % A funnel: any place the set gets smaller -- a cap, a top-k, a filter.
  % Twin: [cap 3 per source].
  pics/cgfunnel/.style={code={
    \path[fill=\cgPicColour!18, draw=\cgPicColour Line, line width=0.8pt]
      (-3.6mm,3.4mm) -- (3.6mm,3.4mm) -- (0.9mm,-0.4mm) -- (0.9mm,-3.8mm)
      -- (-0.9mm,-3.8mm) -- (-0.9mm,-0.4mm) -- cycle;
  }},
  % A chip: anything that computes -- a model, an encoder, a reranker.
  % Twin: [bi-encoder].
  pics/cgchip/.style={code={
    \path[fill=\cgPicColour!15, draw=\cgPicColour Line, line width=0.8pt,
          rounded corners=0.8pt] (-3mm,-3mm) rectangle (3mm,3mm);
    \path[fill=\cgPicColour!45, draw=\cgPicColour Line, line width=0.6pt,
          rounded corners=0.5pt] (-1.4mm,-1.4mm) rectangle (1.4mm,1.4mm);
    \foreach \i in {-1.5,0,1.5}{
      \path[draw=\cgPicColour Line, line width=0.7pt]
        (\i mm,3mm) -- (\i mm,4.3mm)  (\i mm,-3mm) -- (\i mm,-4.3mm)
        (3mm,\i mm) -- (4.3mm,\i mm)  (-3mm,\i mm) -- (-4.3mm,\i mm);}
  }},
  % A shield: the one blocking check, at most once per figure.
  % Twin: [gate].
  pics/cgshield/.style={code={
    \path[fill=\cgPicColour!16, draw=\cgPicColour Line, line width=0.95pt]
      (0,4.2mm) -- (3.4mm,2.6mm) -- (3.4mm,-1mm)
      .. controls (3.4mm,-3mm) and (1.6mm,-3.8mm) .. (0,-4.4mm)
      .. controls (-1.6mm,-3.8mm) and (-3.4mm,-3mm) .. (-3.4mm,-1mm)
      -- (-3.4mm,2.6mm) -- cycle;
    \path[draw=\cgPicColour Line, line width=1pt, line cap=round]
      (-1.5mm,0.2mm) -- (-0.4mm,-1.1mm) -- (1.7mm,1.6mm);
  }},
}%
% >>> end chitragupta figure style <<<
\usetikzlibrary{arrows.meta,positioning,fit,backgrounds,calc,shadows.blur}%
\begin{tikzpicture}[cg, node distance=7mm,
                    pt/.style={text width=33mm, minimum height=14mm}]
  % `pt` sizes a part and sets no role, so it goes after the role: a role
  % style restarts from `cgbox` and would reset a width set before it.
  % ---- panel (a) -----------------------------------------------------------
  \node[cgbox, pt]                     (src-a)   {\cglab{Web search}{live pages}};
  \node[cgbox, pt, below=of src-a]     (draft-a) {\cglab{Draft}{the model writes}};
  \node[cgnone, pt, below=of draft-a]  (check-a) {\cglab{No check}{keys as written}};
  \node[cgbox, pt, below=of check-a]   (out-a)   {\cglab{Publish}{anything cited}};
  % ---- panel (b): its column, each part on its match's row -----------------
  \node[cgbox, pt, right=16mm of src-a] (src-b)  {\cglab{Corpus}{your parsed PDFs}};
  \node[cgbox, pt] (draft-b) at (src-b |- draft-a) {\cglab{Draft}{the model writes}};
  \node[cgkey, pt] (check-b) at (src-b |- check-a) {\cglab{Verify}{every citekey}};
  \node[cgbox, pt] (out-b)   at (src-b |- out-a)   {\cglab{Publish}{only what passed}};
  % ---- the two columns -----------------------------------------------------
  \draw[cgedge] (src-a) -- (draft-a);
  \draw[cgedge] (draft-a) -- (check-a);
  \draw[cgedge] (check-a) -- (out-a);
  \draw[cgedge] (src-b) -- (draft-b);
  \draw[cgedge] (draft-b) -- (check-b);
  \draw[cgedge] (check-b) -- (out-b);
  % ---- cards and letters ---------------------------------------------------
  \begin{scope}[on background layer]
    \node[cgzone=cgMute, fit=(src-a)(draft-a)(check-a)(out-a)] (panel-a) {};
    \node[cgzone=cgMute, fit=(src-b)(draft-b)(check-b)(out-b)] (panel-b) {};
  \end{scope}
  \node[below=2mm of panel-a, font=\normalsize\bfseries] (label-a) {(a) typical writer};
  \node[below=2mm of panel-b, font=\normalsize\bfseries] (label-b) {(b) this pipeline};
\end{tikzpicture}%
